\subsection{测度的无穷乘积}

设 \((\mathbf{X}_{\lambda})_{\lambda \in \mathbf{L}}\) 是紧空间的一个族，且对每个 \(\lambda \in \mathbf{L}\) ，设 \(\mu_{\lambda}\) 是 \(\mathbf{X}_{\lambda}\) 上的一个测度。我们沿用 No.~5 例子中的记号，于是特别地，\(\mu_{\mathrm{J}} = \bigotimes_{\lambda \in \mathrm{J}} \mu_{\lambda}\) 对每个有限子集 \(J\) ⊂ \(L\) 成立。

命题 9. --- 假设诸测度 \(\mu_{\lambda}\) 都是正测度，且族 \((\mu_{\lambda}(1))_{\lambda \in L}\) 在 \(\mathbf{R}_{+}\) 中可乘（乘积可能为 0）。则存在唯一的测度 \(\mu\) 于 X 上，使得对 L 的每个有限子集 J 和每个函数 \(f_{J} \in \mathcal{C}(X_J; C)\) ，

\[
\mu (f _ {\mathrm{J}} \circ \mathrm{pr} _ {\mathrm{J}}) = \mu_ {\mathrm{J}} (f _ {\mathrm{J}}) \prod_ {\lambda \in \mathrm{L} - \mathrm{J}} \mu_ {\lambda} (1).\tag{18}
\]

此外，测度 \(\mu\) 是正测度，其总质量由

\[
\mu (1) = \prod_ {\lambda \in L} \mu_ {\lambda} (1).\tag{19}
\]

设 \(\mathbf{F}\) 是由 \(\mathbf{X}\) 上形如 \(f_{\mathrm{J}} \circ \mathrm{pr}_{\mathrm{J}}\) 的函数所构成的向量空间，其中 \(\mathbf{J}\) 取遍 \(\mathbf{L}\) 的有限子集组成的有向集，而 \(f_{\mathrm{J}} \in \mathcal{C}(\mathrm{X}_{\mathrm{J}}; \mathbf{C})\) ；我们也把 \(\mathbf{F}\) 称为 \(\mathbf{X}\) 上只依赖于有限个变量的连续函数所成的空间。No.~5 的引理 3 表明 \(\mathbf{F}\) 在 \(\mathcal{C}(\mathrm{X}; \mathbf{C})\) 中稠密，这就证明了唯一性断言。 若 \(\mu_{\lambda_0} = 0\) 对某个 \(\lambda_0 \in \mathbf{L}\) 成立，则测度 \(\mu = 0\) 即满足要求，因为在 (18) 的右端中，\(\mu_{\mathrm{J}}(f_{\mathrm{J}}) = 0\) 当 \(\lambda_0 \in \mathbf{J}\) 时成立，而 \(\prod_{\lambda \in \mathbf{L} - \mathbf{J}} \mu_{\lambda}(1) = 0\) 当 \(\lambda_0 \notin \mathbf{J}\) 时成立。于是可以假设 \(\mu_{\lambda} \neq 0\) 对所有 \(\lambda \in \mathbf{J}\) 成立，又因诸测度 \(\mu_{\lambda}\) 是正的，这就蕴涵 \(\mu_{\lambda}(1) \neq 0\) 对所有 \(\lambda \in \mathbf{L}\) 成立。 于是置 \(\mu_{\lambda}' = \mu_{\lambda}/\mu_{\lambda}(1)\) （对每个 \(\lambda \in \mathbf{L}\) ），则 \(\mu_{\lambda}'\) 是 \(X_{\lambda}\) 上满足 \(\mu_{\lambda}'(1) = 1\) 的正测度。从而由 No.~5 的命题 8 可知，存在正测度 \(\mu'\) 于 \(X\) 上，总质量为 1，使得 \(\mu'(f_{\mathrm{J}} \circ \mathrm{pr}_{\mathrm{J}}) = \mu_{\mathrm{J}}'(f_{\mathrm{J}})\) 对每个有限子集 \(J\) ⊂ \(L\) 与每个函数 \(f_{\mathrm{J}} \in \mathcal{C}(\mathrm{X}_{\mathrm{J}}; \mathbf{C})\) 成立。则正测度

\[
\mu = \left(\prod_ {\lambda \in L} \mu_ {\lambda} (1)\right) \mu^ {\prime}
\]

即满足要求，因为

\[
\begin{array}{c} \mu_ {\mathrm{J}} (f _ {\mathrm{J}}) = \mu_ {\mathrm{J}} ^ {\prime} (f _ {\mathrm{J}}) \cdot \prod_ {\lambda \in \mathrm{J}} \mu_ {\lambda} (1), \\ \prod_ {\lambda \in \mathrm{L}} \mu_ {\lambda} (1) = \prod_ {\lambda \in \mathrm{J}} \mu_ {\lambda} (1) \cdot \prod_ {\lambda \in \mathrm{L} - \mathrm{J}} \mu_ {\lambda} (1). \end{array}
\]

命题 9 中定义的测度 \(\mu\) 称为正测度族 \((\mu_{\lambda})_{\lambda \in L}\) 的乘积测度，记作 \(\bigotimes_{\lambda \in L} \mu_{\lambda}\) 。

推论. --- 假设命题 9 的条件成立，且设 \((\mathrm{L}_{\rho})_{\rho\in\mathbb{R}}\) 是 L 的一个分割。则诸测度族 \((\mu_{\lambda})_{\lambda\in\mathrm{L}_{\rho}}\) 中的每一个都有乘积测度，乘积测度族 \(\left(\bigotimes_{\lambda\in\mathrm{L}_{\rho}}\mu_{\lambda}\right)_{\rho\in\mathbb{R}}\) 也有乘积测度，并且

\[
\bigotimes_ {\rho \in \mathrm{R}} \left(\bigotimes_ {\lambda \in \mathrm{L} _ {\rho}} \mu_ {\lambda}\right) = \bigotimes_ {\lambda \in \mathrm{L}} \mu_ {\lambda}.\tag{20}
\]

这可由公式 (18)、(19) 以及 \(R_{+}\) 中可乘族乘积的结合律（GT, IV, §7, No.~5, 注）立即得出。

